\documentclass[10pt,twocolumn,letterpaper]{article}
\usepackage{cvpr}

\usepackage{graphicx}
\usepackage{amsmath,amssymb}
\usepackage{booktabs}
\usepackage[pagebackref,breaklinks,colorlinks]{hyperref}
\usepackage[capitalize]{cleveref}

\begin{document}

\title{Retrieval-Augmented Memory for Long-Horizon Video World Models}

\author{
Aleksandr Trofimov\\
MIT\\
{\tt\small atrof002@mit.edu}
\and
Maryna Bohdan\\
MIT\\
{\tt\small mbohdan@mit.edu}
}

\maketitle

\section{Introduction}
One of the central challenges in robotics is operating under partial
observability: an agent only sees a small part of the environment at any given
moment, but still has to act as if it understands the larger scene. Computer
vision is a key part of solving this problem because visual observations provide
the evidence from which the agent must infer, update, and remember the structure
of the surroundings over time. In many settings, especially embodied and
interactive ones, a system has to reason over time. It must track what has
already been seen, maintain a coherent representation of the environment, and
use that representation to predict future observations. This is difficult
because visual input arrives sequentially, and a useful model therefore needs
not only perception but also memory.

World models are one framework for addressing this problem. Broadly speaking, a
world model learns a compact representation of an environment and uses it to
predict how observations evolve over time. This idea has been influential in
both vision and reinforcement learning. Previous work such as PlaNet and Dreamer
showed that latent dynamics models can support prediction and control directly
from visual observations \cite{hafner2019planet,hafner2020dreamer}. More recent
methods, including VideoGPT and IRIS, demonstrated that latent-token and
transformer-based architectures can model visual sequences effectively
\cite{yan2021videogpt,micheli2023iris}. Together, these results suggest that
predictive latent modeling is generally a promising direction for embodied
visual intelligence.

At the same time, long-horizon consistency remains a clear weakness of current
approaches. A model may generate plausible short-term predictions while still
failing to preserve a stable representation of a scene over longer trajectories.
Once earlier observations fall outside the active context window, the model can
start drifting (object identities may change, spatial layout may become
inconsistent, and previously seen scenes may be replaced by something merely
plausible rather than correct). Transformer-based world models have improved
sequential prediction \cite{chen2022transdreamer}, but finite context still
limits how well they handle information that becomes relevant again only after a
long gap. Recent work such as WorldMem argues that explicit memory may be
necessary for sustaining long-term visual consistency \cite{xiao2025worldmem}.

This issue becomes particularly acute in indoor embodied environments. An agent
may observe a room, move through other parts of the scene, and later return from
a different viewpoint. In that case, a robust world model should still preserve
the identity of the original room: its rough layout, its objects, and their
arrangement. If the model instead alters furniture or predicts a different but
still reasonable-looking room, then it is not really modeling a persistent
environment. This leads to the main question we study in this project: can a
world model maintain scene consistency over time, even when the agent revisits
parts of the environment after a long gap?

We study this question using AI2-THOR and ProcTHOR, which provide controllable
indoor environments, embodied actions, and structured scene metadata
\cite{kolve2017ai2thor,deitke2022procthor}. These simulators are well suited to
our setting because they allow us to generate trajectories and to evaluate
consistency in a controlled way. However, our project is intentionally narrow in
scope, as we are not trying to train a large general-purpose world model across
many homes or room categories. Instead, we focus on a small proof-of-concept
setting, training a compact model centered on just a few rooms, and testing
whether it can preserve one room across revisit. This makes our project feasible
while still isolating the memory problem we care about.

To address this problem, we propose a retrieval-augmented latent world model for
indoor revisit. The model first encodes visual observations into compact latent
representations. A small transformer-based dynamics model then predicts future
latent states from recent observations and actions. We augment this baseline
with an external memory bank that stores selected past latent states together
with simple metadata such as timestep or camera position. During rollout, the
model can retrieve relevant earlier states and condition its predictions on them
in addition to the recent context, and if information about the room has dropped
out of the local context window, targeted retrieval may recover it more
effectively than by relying on the sequence model only.

The project’s contribution is an investigation of whether the model can stay
stable over a long simulation period. First, we estimate long-horizon consistency
through a revisit, which gives a more direct test of memory than standard
short-horizon prediction metrics alone. Second, we examine whether a lightweight
retrieval mechanism can improve consistency in a setting that is realistic under
course-project compute limits. Third, we evaluate the model using metrics that
better suit the question at hand, including visual similarity at revisit time,
object persistence, and coarse agreement with the original room layout.

Because this is an introduction draft, our results are currently prospective
rather than final. We expect retrieval-based memory to matter most as the
temporal gap between first observation and revisit increases. For short
trajectories, recent context may already be sufficient. For longer ones,
however, earlier room information is more likely to be forgotten, and explicit
retrieval should help reduce drift and preserve scene identity more reliably.

Overall, this project studies a small but meaningful version of a broader
computer vision problem: how to maintain a consistent visual model of an
environment over time. By narrowing the task to one-room revisit, we aim to test
whether explicit memory improves long-horizon prediction in a controlled
embodied setting.

\section{AI Use}
Generative AI tools were not used in the preparation of this paper introduction.

\begin{thebibliography}{99}

\bibitem{chen2022transdreamer}
Chang Chen, Yi-Fu Wu, Jaesik Yoon, and Sungjin Ahn.
TransDreamer: Reinforcement Learning with Transformer World Models.
\textit{arXiv preprint arXiv:2202.09481}, 2022.

\bibitem{deitke2022procthor}
Matt Deitke, Eli VanderBilt, Alvaro Herrasti, Luca Weihs, Jordi Salvador, Kiana Ehsani, Winson Han, Eric Kolve, Ali Farhadi, Aniruddha Kembhavi, and Roozbeh Mottaghi.
ProcTHOR: Large-Scale Embodied AI Using Procedural Generation.
\textit{Advances in Neural Information Processing Systems}, 2022.

\bibitem{hafner2020dreamer}
Danijar Hafner, Timothy Lillicrap, Jimmy Ba, and Mohammad Norouzi.
Dream to Control: Learning Behaviors by Latent Imagination.
In \textit{International Conference on Learning Representations}, 2020.

\bibitem{hafner2019planet}
Danijar Hafner, Timothy Lillicrap, Ian Fischer, Ruben Villegas, David Ha, Honglak Lee, and James Davidson.
Learning Latent Dynamics for Planning from Pixels.
In \textit{Proceedings of the 36th International Conference on Machine Learning}, pages 2555--2565, 2019.

\bibitem{kolve2017ai2thor}
Eric Kolve, Roozbeh Mottaghi, Winson Han, Eli VanderBilt, Luca Weihs, Alvaro Herrasti, Matt Deitke, Kiana Ehsani, Daniel Gordon, Yuke Zhu, Aniruddha Kembhavi, Abhinav Gupta, and Ali Farhadi.
AI2-THOR: An Interactive 3D Environment for Visual AI.
\textit{arXiv preprint arXiv:1712.05474}, 2017.

\bibitem{micheli2023iris}
Vincent Micheli, Eloi Alonso, and Fran\c{c}ois Fleuret.
Transformers are Sample-Efficient World Models.
In \textit{International Conference on Learning Representations}, 2023.

\bibitem{xiao2025worldmem}
Zeqi Xiao, Yushi Lan, Yifan Zhou, Wenqi Ouyang, Shuai Yang, Yanhong Zeng, and Xingang Pan.
WORLDMEM: Long-term Consistent World Simulation with Memory.
\textit{arXiv preprint arXiv:2504.12369}, 2025.

\bibitem{yan2021videogpt}
Wilson Yan, Yunzhi Zhang, Pieter Abbeel, and Aravind Srinivas.
VideoGPT: Video Generation using VQ-VAE and Transformers.
\textit{arXiv preprint arXiv:2104.10157}, 2021.

\end{thebibliography}

\end{document}